\hyphenation{ALTA} % to avoid breaking ALTA over two lines

Welcome to the \textbf{24th Annual Workshop of the Australasian Language Technology Association (ALTA 2026)}. ALTA 2026 was held at Deakin University's Downtown campus in Melbourne from 30th November to 2nd December 2026, on the lands of the Wurundjeri Woi-wurrung and Bunurong / Boon Wurrung peoples of the Kulin Nation. We acknowledge them as the Traditional Owners of these lands and pay our respects to their Elders, past and present, as well as to Aboriginal and Torres Strait Islander Elders from other communities who joined us. The workshop brought together researchers, industry experts, and practitioners working in natural language processing (NLP), computational linguistics, and related fields. It provided a forum for exchanging ideas, exploring innovative approaches, and discussing the latest developments in language technology. \newline


The ALTA 2026 program reflected a wide variety of topics relevant to contemporary research in these fields. As per its tradition, ALTA offered two submission options: archival papers (long and short) and non-archival papers (various lengths). We received 56 active archival submissions (excluding desk-rejected and withdrawn papers), of which 23 were accepted following a double-blind peer-review process involving at least two reviewers per paper. This represents an overall acceptance rate of 41.07\%. Of the 56 active submissions, 57.14\% (32) originated from Australia, while 42.86\% (24) came from other countries. The accepted papers represented six countries, with contributions from Australia (17), the United States (2), Sri Lanka (1), Indonesia (1), Egypt (1), and Bangladesh (1). These contributions demonstrate the international reach of ALTA 2026 and its growing engagement with the global NLP research community. \newline


The accepted archival papers cover a broad range of emerging challenges and opportunities in NLP, with a particular focus on the capabilities, reliability, robustness, and responsible use of large language models (LLMs). Key themes include reasoning, hallucination detection, bias, fairness, trust, and consistency in LLMs, alongside advances in retrieval-augmented generation, language agents, and efficient model adaptation. Several papers explore applications in high-stakes domains such as healthcare, education, and legal analysis. The accepted papers also highlight multilingual and low-resource language technologies, including work on Sinhala, Arabic, and Kupang Malay. These contributions reflect the diversity of contemporary NLP, combining methodological advances with practical applications to support more reliable, efficient, inclusive, and trustworthy language technologies. \newline


The non-archival track welcomed submissions reporting on early-stage research, contributions from industry practitioners, and papers previously published at international conferences, including ACL-affiliated conferences and workshops. The track aimed to provide a local forum for disseminating internationally recognised research and fostering discussion among researchers in the region. We received 45 active submissions (excluding withdrawn papers) in the non-archival track, of which 19 were accepted, resulting in an acceptance rate of 42.22\%. Each non-archival submission was reviewed by a member of the workshop's organising committee, comprising the general and program chairs. The accepted non-archival papers were presented as posters. \newline


The workshop featured two tutorials. The morning tutorial, titled 'Transfer Learning in Clinical NLP', was presented by Distinguished Prof. Karin Verspoor, Dr. David Martinez Iraola, and Ms. Wei Han from RMIT University. The afternoon tutorial, titled 'Evaluating AI Agents with LLM-as-a-Judge', was presented by Dr. Parvin Zarei Eskikand and Dr. Terrence Szymanski from SEEK, Ms. Sara Mirabi from Deakin University, and Dr. Yelaman Abdullin from Macquarie University. \newline


This year's panel discussion focused on 'From Chatbots to Teammates: Enabling Effective Human-Agent Collaboration'. The panel brought together Dr. Rai Tumuluri (Openstream.ai), Associate Prof. Zoie Wilkins-Wong (Swinburne University), and Dr. Johanne Trippas (RMIT University). We thank the panelists for sharing their perspectives and insights on this important topic. We also acknowledge the contributions of the shared task chair, Dr. Diego Molla-Aliod, from Macquarie University, and all participants. The shared task participants presented their work in short talks during the workshop, providing further opportunities for discussion and knowledge exchange. \newline


We are pleased to acknowledge our distinguished keynote speakers, Prof. Eduard Hovy (University of Melbourne), Distinguished Prof. Karin Verspoor (RMIT University), and Distinguished Prof. Mohit Bansal (UNC Chapel Hill), for sharing their insights and expertise with the ALTA 2026 community. We also wish to express our sincere appreciation to Prof. Massimo Piccardi (University of Technology Sydney), Dr. Ming Liu (Deakin University), and Dr. Mohammed Reda Bouadjenek (Deakin University) for generously serving as judges for the Best Paper Awards. \newline


We would like to express our sincere gratitude to the members of the organising committee and the volunteers for their dedication and hard work in preparing for and running ALTA 2026. We also extend our heartfelt appreciation to all reviewers for their time, expertise, and effort in providing thoughtful, detailed, and constructive feedback on the submissions. Their contributions have been invaluable to the success and quality of the workshop. \newline


ALTA 2026 gratefully acknowledges the generous support of its sponsors: Defence Science and Technology Group and Openstream.ai (Platinum); Google, MongoDB, and SEEK (Silver); and the University of Melbourne (Bronze). We also extend our sincere thanks to Deakin University for hosting ALTA 2026 at its Downtown campus in Melbourne. \newline


The success of ALTA 2026 would not have been possible without the invaluable contributions of all those involved. We thank everyone who joined us at Deakin University and contributed to making the workshop an engaging and rewarding experience. We hope that the workshop provided an inspiring opportunity for the exchange of ideas and perspectives and helped strengthen connections within the international NLP research community. \newline


\hspace*{\fill}Bahadorreza Ofoghi\newline
\hspace*{\fill}\textit{ALTA 2026 General Chair}\newline\newline
\hspace*{\fill}Massimo Piccardi\newline
\hspace*{\fill}Ming Liu\newline
\hspace*{\fill}Mohamed Reda Bouadjenek\newline
\hspace*{\fill}\textit{ALTA 2026 Program Chairs}